%  LaTeX support: latex@mdpi.com 
%  For support, please attach all files needed for compiling as well as the log file, and specify your operating system, LaTeX version, and LaTeX editor.

%=================================================================
\documentclass[jimaging,article,submit,pdftex,moreauthors]{Definitions/mdpi} 

%=================================================================

% MDPI internal commands
\firstpage{1} 
\makeatletter 
\setcounter{page}{\@firstpage} 
\makeatother
\pubvolume{1}
\issuenum{1}
\articlenumber{0}
\pubyear{2023}
\copyrightyear{2023}
%\externaleditor{Academic Editor: Firstname Lastname}
\datereceived{23 February 2023} 
\daterevised{2 April 2023} 
\dateaccepted{} 
\datepublished{} 
\hreflink{https://doi.org/} % If needed use \linebreak
%\doinum{}

%=================================================================
% Add packages and commands here. The following packages are loaded in our class file: fontenc, inputenc, calc, indentfirst, fancyhdr, graphicx, epstopdf, lastpage, ifthen, lineno, float, amsmath, setspace, enumitem, mathpazo, booktabs, titlesec, etoolbox, tabto, xcolor, soul, multirow, microtype, tikz, totcount, changepage, attrib, upgreek, cleveref, amsthm, hyphenat, natbib, hyperref, footmisc, url, geometry, newfloat, caption

\graphicspath{{figures/}} % path to folder with figures
\usepackage{subcaption}
\usepackage{listings}
\usepackage{xcolor}
\usepackage{gensymb} % for \textdegree
\usepackage{tabularx}
\usepackage{array}
\usepackage{makecell}
\usepackage{multirow}

\usepackage{listings}
\usepackage{xcolor}

\definecolor{deepblue}{rgb}{0,0,0.5}
\definecolor{deepred}{rgb}{0.6,0,0}
\definecolor{deepmagenta}{RGB}{195,12,214}
\definecolor{deepgreen}{RGB}{29,122,30}
\definecolor{mintgreen}{RGB}{192,249,192}
\definecolor{codepurple}{RGB}{204, 0, 153}
\definecolor{LightCyan}{rgb}{0.88,1,1}
\definecolor{GainsboroPurple}{RGB}{247,176,247}
\definecolor{LightSlateGrey}{RGB}{124,118,118}
%    
\lstdefinestyle{mystyle}{
    commentstyle=\color{LightSlateGrey},
    keywordstyle=\color{deepgreen}\bfseries,
    emphstyle=\color{deepred},
    numberstyle=\tiny\color{deepmagenta},
    stringstyle=\color{codepurple},
    basicstyle=\ttfamily\scriptsize,
    breakatwhitespace=false,         
    breaklines=true,                 
    captionpos=b,                    
    keepspaces=true,                 
    numbers=left,                    
    numbersep=5pt,                  
    showspaces=false,                
    showstringspaces=false,
    showtabs=false,                  
    tabsize=2
}
\lstset{style=mystyle}

%================================================================= in Khartoum Region of Sudan, NE Africa
% Full title of the paper (Capitalized)
\Title{Multispectral Satellite Image Analysis for Computing Vegetation Indices by R in the Khartoum Region of Sudan, Northeast Africa}

% MDPI internal command: Title for citation in the left column
\TitleCitation{Multispectral Satellite Image Analysis for Computing Vegetation Indices by R in the Khartoum Region of Sudan, Northeast Africa}

% Author Orchid ID: enter ID or remove command
\newcommand{\orcidauthorA}{0000-0002-5759-1089} % Polina Add \orcidA{} behind the author's name
\newcommand{\orcidauthorB}{0000-0002-6461-1551} % Olivier Add \orcidB{} behind the author's name

% Authors, for the paper (add full first names)
\Author{Polina Lemenkova $^{1}$*\orcidA{} and Olivier Debeir $^{1}$\orcidB{}}

%\longauthorlist{yes}

% MDPI internal command: Authors, for metadata in PDF
\AuthorNames{Polina Lemenkova and Olivier Debeir}

% MDPI internal command: Authors, for citation in the left column
\AuthorCitation{Lemenkova, P.; Debeir, O.}
% If this is a Chicago style journal: Lastname, Firstname, Firstname Lastname, and Firstname Lastname.

% Affiliations / Addresses (Add [1] after \address if there is only one affiliation.)
\address{%
$^{1}$ \quad Laboratory of Image Synthesis and Analysis (LISA), École polytechnique de Bruxelles (Brussels Faculty of Engineering), Université Libre de Bruxelles (ULB). Building L, Campus du Solbosch, ULB -- LISA CP165/57, Avenue Franklin D. Roosevelt 50, Brussels 1050, Belgium.
}

% Contact information of the corresponding author
\corres{Correspondence: polina.lemenkova@ulb.be; Tel.: +32 471 86 04 59 (P.L.)}

% Current address and/or shared authorship
%\firstnote{Current address: Affiliation 3.} 

\abstract{Desertification is one of the most destructive climate-related issues in the Sudan-Sahel region of Africa. As the assessment of desertification is possible by satellite image analysis using Vegetation Indices (VIs), this study reports on the technical advantages and capabilities of scripting the 'raster' and 'terra' R-language packages for computing the VIs. The test area which was considered includes the region of the confluence of the Blue Nile and White Nile in Khartoum, southern Sudan, Northeast Africa and the Landsat 8-9 OLI/TIRS images taken for the periods of 2013, 2018 and 2022 were chosen as test datasets. The VIs used here are robust indicators of plant greenness, and combined with vegetation coverage are essential parameters for environmental analytics. Five VIs were calculated to compare both the status and dynamics of vegetation through the differences between the images collected within the nine year span. Using scripts for computing and visualising the VIs over Sudan demonstrates previously unreported patterns of vegetation to reveal climate--vegetation relationships. The ability of the R packages 'raster' and 'terra' to process spatial data was enhanced through scripting to automate image analysis and mapping , and choosing Sudan for the case study enables us to present new perspectives for image processing.
}
% Keywords
\keyword{Remote sensing; Sudan; Khartoum; Nile River; image analysis; image processing; information retrieval; programming; Africa; computer vision} 

\PACS{91.10.Da; 91.10.Jf; 91.10.Sp; 91.10.Xa; 96.25.Vt; 91.10.Fc; 95.40.+s; 95.75.Qr; 95.75.Rs; 42.68.Wt}
\MSC{86A30; 86-08; 86A99; 86A04}
\JEL{Y91; Q20; Q24; Q23; Q3; Q01; R11; O44; O13; Q5; Q51; Q55; N57; C6; C61}

%%%%%%%%%%%%%%%%%%%%%%%%%%%%%%%%%%%%%%%%%%
\begin{document}

\section{Introduction}

\subsection{Background}

In this paper, we introduce the application of R libraries for remote sensing data processing in the context of computing vegetation indices (VIs) that indicate desertification processes. Further, we focus on calculating several satellite-derived VIs by discussing their technical details and differences. The test dataset includes the Landsat 8-9 OLI/TIRS images collected on Khartoum city around the confluence of the White Nile and the Blue Nile. The images were used to analyse the changes in the vegetation patterns and greenness through analysing the time series and visualising the dynamics in vegetation coverage. The standard methods for such tasks are based on using the traditional GIS software \cite{AHMED201829} and involve many\hl{ }routine operations \cite{Hawash,YOUSSEF2020103777,6621869,Aldoma}. Instead of using such conventional approaches, we apply several libraries of R such as 'raster' and 'terra' and auxiliary packages, in a principled way designed to obtain a better workflow for the image analysis involved.

The Earth observing satellite sensors providing remote sensing data enable to highlight the links between the hydrological and climate processes leading to desertification and vegetation responses in the environmental studies. The importance of the use of the remote sensing data and advanced methods of image processing is raised in earlier studies\cite{su14137871,ELBASTAWESY2017252,SOSNOWSKI201651,GORSEVSKI201364}. Satellite images can be used to analyse the links between the complex hydrological and climate processes and vegetation responses that lead to desertification. For example, Landsat images are known to be a reliable source of data for relatively accurate techniques for classifying time-series and detecting forest and land cover types\cite{WILSON2002385,Lemenkova202227,4293065,SALIH2017S21}, computing vegetation indices \cite{5416935,Lemenkova202229,1026428} and specifically desertification \cite{RIVERAMARIN2022104829} to show the advance or retreat of arid areas using analysis of satellite images. Specifically, for the Sudan and Nile River Basin, the Landsat data are the precious source of information, due to data scarcity \cite{Ali} such as regular measurements of rainfall, streams run-off, and  weather stations data. This highlights the importance of the remote sensing data for climate-related studies of Sudan.

GIS was one of the earliest computer-based application for satellite image processing developed for the empirical investigations of the environmental parameters derived from the remote sensing data. Since the 1990s, researchers were able to accurately compute the vegetation indices derived from the images using GIS. Kumar et al. have studied performance of the ERDAS Imagine in computing vegetation indices to assess desertification in semi-arid regions, emphasising that the NDVI classification is suitable for assessing land degradation using the NDVI pixel range \cite{KUMAR2022100578}. Rashid et al. have examined the links between the NDVI and the salinity characteristics in the surface water, groundwater, and soils using the ArcGIS \cite{RASHID2023100314}. Further, Ezaidi et al. applied ENVI GIS for computing the NDVI for a time series of the Landsat images to assess the degradation in vegetation \cite{EZAIDI2022100800} All these studies illustrate the traditional approaches for computing vegetation indices by GIS.

As a general overview, GIS enables to calculate vegetation indices through the manipulation of Landsat images using map algebra. However, the efficacy of the computational approaches must be questioned, as more advanced tools of spatial data analysis enable the application of modelling methods for computing the VIs. For instance, \cite{MARTINEZ2023115155} demonstrated the use of by Generalised Additive Models to study the relationships between vegetation types and NDVI. Also, \cite{OMAR2021e01020} demonstrates the use of Holt-Winters modelling to detect trend and seasonal vegetation pattern using Python's Scikit-Learn library, and \cite{WANG2022102704} demonstrates the use of partial least squared regression (PLSR) model in the SIMCA-P software to track the non-linear relationships between the NDVI and climatic variables. Recent development of the machine learning tools through the use of the Python programming language have enabled the use of Convolutional Neural Network (CNN) and Deep Neural Networks (DNN) for environmental analysis and estimation of the VIs \cite{BARRIGUINHA2022103069,ROY2021100582,LI2022102818,DAVIDSON2022107396}. 

Compared to these and similar programming approaches, the use of R in the context of the computing of VIs presents a more straightforward approach. Similarly to GIS approaches, R is adjusted to the processing of the Landsat bands; at the same time, the flexibility of R being programming language allows full control of the data processing and a high level of automation. Examples of R used for computing VIs include the 'raster' package for computing the NDVI \cite{LEVEAU2021127199}, the 'npphen' R package to assess the NDVI anomalies and evaluating spatio-temporal changes \cite{CHAVEZ2023103138}, the 'akima' R package to interpolate environmental parameters \cite{KLIMAVICIUS2023171} or the 'RStoolbox' package for computing the NDVI \cite{PRAVALIE2022108629}. Such methods are comparable to our approach, however, they are largely tied to specific problem formulation applied to the target areas. To the best of our knowledge, no previous studies reported on the use of R for computing the VIs in Sudan. This study fills this gap and highlights the benefits of using advanced techniques to compute the VIs specifically for Sahelian landscapes of southern Sudan, prone to desertification.

\subsection{\textcolor{blue}{Objectives and Motivation}}

\hl{The main purpose of the current paper is to apply an R-based method of satellite image processing that can extract information on several vegetation indices for evaluating the vegetation patterns in arid environmental setting of southern Sudan. The method of R utilizes the concept of scripting libraries and suggests a way to obtain values of spectral bands for computing diverse vegetation indices that describe the vegetation greenness and health. Subsequently is builds an approach to analyse remote sensing data associated with spectral signature of land surface features visible on the space borne data. The resulting series of the computed vegetation indices serve can be evaluated to support environmental monitoring and mapping in semi-arid and arid areas of Eastern Africa.}

\hl{In this paper, we address the issue of using R programming language in geospatial studies with a particular case of satellite image processing. Compared with the existing literature, the main contributions of this paper are highlighted as follows:}
\begin{itemize}
	\item \hl{To the best of our knowledge, this is the first paper in the literature that focuses on the application of R libraries for computing several vegetation indices over the area of Khartoum, Sudan;}
	\item \hl{As opposed to previous studies, we present the application of 'raster' and 'terra' packages of R for remote sensing data analysis instead of the traditional GIS software;}
	\item \hl{We further extend the use of Landsat 8-9 OLI/TIRS sensors to extract geospatial information from multi-spectral images;}
	\item \hl{We evaluate comprehensively the performance of different vegetation indices for the case of semi-arid vegetation patterns in southern Sudan. To this end, we apply the R algorithms to calculate five different VI which differ in computational approaches and tuned to diverse environmental aspects of vegetation: 1) Normalized Difference Vegetation Index (NDVI), 2) Normalized Difference Water Index (NDWI), 3) Infrared Percentage Vegetation Index (IPVI), 4) Optimized Soil Adjusted Vegetation Index (OSAVI), and 5) Green Normalized Difference Vegetation Index (GNDVI);}
	\item \hl{We provide an overview of major environmental aspects of Sudan which include recent issues of desertification related to semi-arid climate, land cover changes and hydrology of the Nile which controls the cycle of vegetation growth in Sudan}.
	\item \hl{A summary of R scripts used for computing and mapping vegetation indices is reported, to provide the reader with the technical reference of the applied methods.}
\end{itemize}

\hl{The remainder of this paper is organised as follows. We review the environmental setting of the study area with related works in the following subsection. In Section 2. Materials and Methods we provide a description of the data and methods used in this study. We describe the data used and report their major characteristics in subsection 2.1. Data We introduce an R language concept, present a principle of R-base libraries for computing vegetation indices and propose the use of 'raster' and 'terra' libraries in subsection 2.2. Methodology. In subsection 2.3. Calculating Vegetation Indices, we provide a comprehensive experimental evaluation of computing several vegetation indices with provided details on the algorithms,  differences in their implementation and target applications. In Section 3. Results, we report on the results of five computed vegetation indices and compare their performance on the vegetation of southern Sudan. We conclude this paper with a discussion in Section 4 and provide closing remarks in Section 6. The main metadata of the satellite images and programming listings of the R scripts are provided in the Appendices.}

\subsection{Study Area}

Desertification is one of the most destructive climate-related issue in Sudan which is located southwards of the Sahara Desert in the Sudan-Sahel region of Africa, (Figure \ref{fig01}). The desertification results from various factors including climate change and human activities. The climate drivers include rising temperatures, decreased relative humidity \cite{Alvi}, increased soil salinization \cite{DAWELBAIT201245}, seasonality and decline of rainfall \cite{MOHAMED20223265}, instability and irregularity of wet and dry seasons \cite{HULME198689}, variability of precipitation \cite{Elagib,Ayoub} and overall water deficit \cite{akhtar1994methods}. Recent studies such as \cite{SALIH2017S31} reported frequent and severe drought periods during the last decades resulting in several cycles of famine in Sudan with affected farmers and population. For instance, chronic drought was reported in the Sahel region of Darfur \cite{FULLER198739}, droughts in central Sudan \cite{Salih2022}, or eastern Sudan \cite{TARIKU2021144863,SULIEMAN2012132}. 

Other factors include changed patterns of cloudiness and atmospheric radiation which results in high evapotranspiration in the extreme north near Sahara, and southern Sahel and savannah and affects vegetation growth \cite{ElKarouri1986}. \hl{Hence, the environment challenges in the semi-arid climate of Sudan affect vegetation patterns which are controlled by the regional climate setting and variations in rainfall, evaporation, cloudiness, atmospheric radiation and other climate-related processes. Recent droughts and decrease in rainfall affected vegetation growth and health due to the lack of humidity in the air. As a consequence, this initiated gradual processes of desertification in Saharan-Sahelian Africa, expansion of drylands, degradation of the landscapes.}

\begin{figure}[H]
\centering
	\includegraphics[width=14.0 cm]{Fig_01.jpg}
	\caption{Topographic map of Sudan. Mapping software: Generic Mapping Tools (GMT) scripting toolset version 6.1.1, \cite{Wessel}. Data source: GEBCO/SRTM. Cartography source: authors.
	\label{fig01}}
\end{figure}

These processes trigger environmental degradation, especially in vulnerable zones of central and southern Sahelian zone of Sudan \cite{Pearce,YANG2022106213} and savannah \cite{Falkenmark}. The most endangered region of Sudan prone to the risks of desertification lies between 12° to 18°N, near its capital -- Khartoum. The flat topography of the country exacerbates the desertification problems and increases the effects from dust and sand storms that often occur in the northern regions of Sudan \cite{khiry2006mapping,9635812,7948717,7381366}. Specifically, this includes the effects from the creeping sands and dunes from the areas of Sahara, Nubian and Libyan deserts, Figure \ref{fig01}. The  consequence os this is deforestation, soil desiccation, changed land cover types \cite{Biro} and decreased biodiversity. The latter includes the extinction of wild species in oases, and consequences of land cover changes include the imbalance of ecosystems, spread of invasive species \cite{insects10110405} and disturbance of natural habitat patterns \cite{Fuller}. 

70\% of the territory of Sudan belongs to the Nile basin \cite{Melesse} with 5 out of the the 6 Cataracts of the Nile located within Sudan (Figure \ref{fig01}). Moreover, roughly 2/3 of the territory of the Nile lies within Sudan with its major tributaries -- the While Nile and Blue Nile -- joining near Khartoum \cite{Hamad}. Hence, the specifics of the Khartoum region are tightly connected to the Nile River system and the confluence of the White Nile and Blue Nile, Figure \ref{fig02}. The impact the Nile River has on the environment includes occasional floods, soil erosion and succession of vegetation on the islands and banks \cite{Halwagy}, accumulation of sedimentation \cite{BILLI201012}, or growth of aquatic weeds \cite{PACINI2016242}. While providing a specific place for wetland habitat of rare species, the Nile river is also essential an water artery of the country, and plays multiple crucial roles. It is a source of hydropower in Sudan \cite{AYIK20232229}, an important transport waterway, and provides water resources for 85\% of the population and irrigated agriculture \cite{Kauetal2023}. \hl{Moreover, the functionality and structure of the riparian vegetation in Sudan strongly depend upon the hydrology of Nile, especially in its southern region near the confluence of the White Nile and the Blue Nile.}

\begin{figure}[H]
\makebox[1.0\linewidth][r]{%
	\begin{subfigure}[b]{.5\textwidth}
		\centering
			\includegraphics[width=7cm]{Fig_02a.jpg}
		\caption{}
	\end{subfigure}%
	\begin{subfigure}[b]{.5\textwidth}
		\centering
			\includegraphics[width=7cm]{Fig_02b.jpg}
		\caption{}
	\end{subfigure}%
}
\caption{The confluence of the White Nile and Blue Nile in Khartoum, visualised based on the map and aerial images from the Google Earth: (\textbf{a}) Google Earth Map, (\textbf{b}) Google Earth Image.}\label{fig02}
\end{figure}

Social factors of desertification include processes related to land management and agriculture policymaking \cite{KOBAYASHI2020257}. Although agriculture and livestock play an important role in the economy of Sudan and major sources of livelihood \cite{sulieman2010expansion}, unsustainable agricultural expansion and farming contribute to Sudan's desertification \cite{Ibrahim} since uncontrolled agriculture results in soil erosion, compaction and nutrient depletion, changed irrigation patterns \cite{AHMED2020106064} and the increase in pesticides \cite{Nesser} which contribute to land cover changes. Moreover, cattle overgrazing also modifies the landscape structure and contributes to desertification \cite{Ahmed,Sulieman,Akhtar}. The intense practice of urban agriculture in Khartoum aims at meeting the increasing food demands of the growing population \cite{SCHUMACHER2009399} which results in the additional pressure on the environment on the one hand, and gradual increase of the built-up area on the other. The confluence of the Blue Nile and White Nile forming its major artery of the Greater Nile near Khartoum results in a high concentration of population in this district. Furthermore, consequences of the desertification also drive people into Khartoum, thereby increasing the urbanisation and the environmental pressure in the region \cite{Babiker1982,Mohammed}, and the Atbara River (Red Nile) has high mining activities related to resource exploration \cite{SASMAZ2020106539}. 

The desertification leads to many negative environmental consequences such as land degradation in arid and semi-arid regions of Sudan \cite{Khiry,ELNASHAR2022152925}. Social consequences of desertification in Sudan include increased droughts and associated famine due to affected food production and soil degradation in the Sahelian zone \cite{Subramaniam}. The environmental changes in the past highlighted links between the human occupation and responses from environment \cite{HONEGGER2015141}. For instance, deep ploughing of the surface increased the susceptibility of soil to wind erosion, which leads to a severe decline in fertility and formation of sand dunes \cite{Abdi}. Other examples of the effects of the desertification include reduced biological productivity \cite{KHALIFA2018790,Jahelnabi} and decrease in plant biomass, increase of the extremely dry desert areas. The cumulative effect of these factors include the unpredictable and severe droughts, desiccation or aridification and dryland degradation or desertification in the dryland regions \cite{Darkoh}. 

Climatic trends are reflected in the hydrology of the Nile which shows high levels of the variability in river flow records \cite{CONWAY200599} and variations in peak discharge \cite{Babiker}. For instance, the increasing trend in natural water storage variation of the Nile can be detected in northern Sudan \cite{ABDELMALIK2019103596}. Such fluctuations in the hydrology of Nile necessarily affect the vegetation coverage in the banks and surrounding areas. Besides, the patterns of the vegetation in the Khartoum Province are deeply connected to the regional geologic setting \cite{ELSHEIKH201182}, the geomorphic structure and dominating soil types that control the distribution of the specific plants \cite{Mahmoud}. The White Nile is sensitive to the rainfall and evaporation in lakes and swamps due to the specific hydrology of the region. Therefore, its response to occasional climatic events can be measured over long timescales \cite{Sene} and the analysis of the vegetation patterns and soil erosion can be used for detecting vulnerable regions prone to desertification \cite{Bakhit}. Nevertheless, despite efforts being put on water resources conservation and rational exploitation of resources \cite{MOHAMEDALI2003237,w10111579}, the environmental sustainability in Sudan region remains a challenge.


 
%%%%%%%%%%%%%%%%%%%%%%%%%%%%%%%%%%%%%%%%%%

\section{Materials and Methods}

%----------- subsection ----------->
\subsection{Data}

The satellite images were selected from the United States Geological Survey (USGS) EarthExplorer repository (\href{https://earthexplorer.usgs.gov/}{https://earthexplorer.usgs.gov/}). The choice of the data is explained by the availability of the Landsat OLI/TIRS scenes, sensor spectral band pass, high spatial resolution of 30 m across the multispectral channels (1 to 7), open availability of the cloud-free images (0\% of cloudiness) and regular survey which provides the comparable gap between the target years (2013 and 2022) where the images were collected for analysis of changes in vegetation (Figure \ref{fig03}).

\begin{figure}[H]
\makebox[1.0\linewidth][r]{%
	\begin{subfigure}[b]{.40\textwidth}
		\centering
			\includegraphics[width=0.90\textwidth]{Fig_03a.jpg}
		\caption{2013}
	\end{subfigure}%
	\begin{subfigure}[b]{.40\textwidth}
		\centering
			\includegraphics[width=0.90\textwidth]{Fig_03b.jpg}
		\caption{2018}
	\end{subfigure}%
	\begin{subfigure}[b]{.40\textwidth}
		\centering
			\includegraphics[width=0.90\textwidth]{Fig_03c.jpg}
		\caption{2022}
	\end{subfigure}%
}
\caption{Landsat 8 OLI/TIRS image of White and Blue Nile in Khartoum, Sudan, in natural RGB colours in December: (\textbf{a}) 20.12.2013, (\textbf{b}) 18.12.2018, (\textbf{c}) 21.12.2022}\label{fig03}
\end{figure}

Among the available Landsat produces, the Landsat OLI/TIRS sensor was selected, as it has better technical and spectral characteristics compared to older sensors. Thus, compared to the Landsat-7 ETM+, the Landsat 8-9 OLI/TIRS has narrower spectral bands, higher calibration, finer radiometric resolution and geometry, better signal to noise parameters \cite{ROY201657}. This results in a more narrow and precise bandwidth in red and NIR channels which differs it from the older sensors. For instance, the Landsat 8-9 OLI/TIRS sensor has 0.64-0.67 \textmu m in the Red band (Band 4), and 0.85-0.88 \textmu m in the NIR (Band 5), while the Landsat 7 ETM+ sensor has coarser range of 0.63-0.69 \textmu m in the Red (Band 3) and 0.77-0.90 \textmu m in NIR (Band 4). The spectral characteristics of the older Landsat products (Landsat 4-5 TM and Landsat 1-5 MSS) has even coarser parameters. Since the Landsat 8 OLI/TIRS was launched on February 11, 2013, there is no earlier images, and the earliest possible date for the cloud-free image on target month of December was 20/12/2013. The images were taken within the 5 year gap between 12/2013 and 12/2018, and 4 year gap between 12/2018 and 12/2022 which gives the comparable intervals of the periods between the scenes.

The summary of the data is presented in Table \ref{tab1}. All the images have 0\% cloudiness and Path/Row ID 173/49. The Sensor ID is common for all the scenes: Landsat 8-9 OLI/TIRS (Operational Land Imager and Thermal Infrared Sensor), Collection 2 Level-2. Image source: the USGS \cite{EROS}. The major technical and geodetic metadata common for all the images are the following: Datum and Ellipsoid WGS84; Product Map Projection L1 -- Universal Transverse Mercator (UTM); UTM Zone 36; Sensor ID OLI TIRS; Landsat Processing Software Version LPGS\_16.2.0. The full list of metadata is provided in Table \ref{tabApp}.

\begin{table}[H]
\scriptsize
\caption{Satellite images used for computing the VIs: Landsat-8 OLI/TIRS collected from the USGS \textsuperscript{1}.\label{tab1}}
%	\begin{adjustwidth}{-\extralength}{0.5cm}
%		\newcolumntype{C}{>{\centering\arraybackslash}X}
		\begin{tabularx}{\linewidth}{ p{1.6cm} | p{1.3cm} | p{6.6cm} | p{2.2cm} }
			\toprule
			\textbf{Date} & \textbf{Spacecraft} & \textbf{Landsat Product ID} & \textbf{Scene ID} \\
			\midrule
			2013/12/20 & Lands. 8 & \makecell{LC08\_L2SP\_173049\_20131220\_20200912\_02\_T1} & LC81730492013354LGN01 \\
			2018/12/18 & Lands. 8 & \makecell{LC08\_L2SP\_173049\_20181218\_20200830\_02\_T1} & LC81730492018352LGN00 \\
			2022/12/21 & Lands. 9 & \makecell{LC09\_L2SP\_173049\_20221221\_20221224\_02\_T1} & LC91730492022355LGN00 \\		
			\bottomrule
		\end{tabularx}
%	\end{adjustwidth}
%	\noindent{\footnotesize{\textsuperscript{1} }}
\end{table}

%----------- subsection ----------->
\subsection{Methodology}

With this study, we process the satellite images using R version 4.2.2 (R Core Team, 2020) and its packages for the computation of VIs and spatial data analysis \cite{RCoreTeam}. The extraction of the VIs for analyses of vegetation health by traditional software, which is conceptually straightforward due to raster algebra, may lead nevertheless to a consequent computing time. The VIs are used as robust and reliable indicators of plant health and greenness because they indicate at the level of leaf chlorophyll in plants. In such a way, the VIs reflect the functional features of canopy and parameters of vegetation coverage essential for environmental and climate analytics.

Over the course of the past 20 years, there has been a significant progress in programming applications in environmental studies towards automation methods in image processing, such as R scripts. Scripting methods in geoinformatics have undergone a strong evolution during the last decades and now they generally result in a higher accuracy, diverse computational approaches in a variety of libraries and pre-programmed packages for data processing. Such technical performance is beneficial for image processing and spatial data analysis compared to traditional GIS methods \cite{6049966,Lemenkova202301,6910592,Lemenkova202228,9323708,1027236}. Considering such advantages of the programming languages in automation of geospatial data analysis, the algorithms of R present a faster and more effective approach to implement image processing scripts. 

In this study, we compute the VIs and analyse the maps based on these indices for the years between 2013 and 2022. In a first step, the necessary libraries 'terra' version 1.7-18, 'RColorBrewer' version 1.1-3 \cite{Neuwirth}, 'Hmisc' version 5.0-1 \cite{Hmisc} and 'pals' version 1.7 \cite{Wright} were activated and a working folder is defined. The 'terra' library is the key package which was used for raster algebra and calculation of the VIs, while the other libraries are used as auxiliary packages for graphical refinements and proper visualisation of the plots. Then, we defined separate functions for each of the VIs using their relevant formulae. Finally, we uploaded the necessary Landsat bands in the active folder and performed calculation of each of the VIs by R scripts, as presented below for each corresponding index.

To compute the VIs, we used the 'terra' R package approach which allows assessment of the vegetation patterns using computational formulae to study vegetation greenness and to assess the increase of desert areas. There are two important advantages of the R approach with regards to estimation of vegetation health using histograms showing data distribution for analysis of the phenological changes in canopies: i) the embedded map algebra is adjusted to describe any type of VIs using function ($vi <- function(img, k, i)$) (see the R scripts in Appendix) that enables to vary Landsat bands for computational purposes in different formulae; and ii) it calculates the data frequency distribution showing the number of pixels that correspond to higher greenness, from which the trends to desertification can be assessed using image analysis for various years. The full scripts used for plotting the Figures \ref{fig01}, \ref{fig04}, \ref{fig05}, \ref{fig06}, \ref{fig07} and \ref{fig08} are provided in the GitHub repository: \href{https://github.com/paulinelemenkova/Mapping_Sudan_Scripts}{https://github.com/paulinelemenkova/Mapping\_Sudan\_Scripts}

%----------- subsection ----------->
\subsection{Calculating Vegetation Indices}

\subsubsection{Normalized Difference Vegetation Index (NDVI)}
The Normalized Difference Vegetation Index (NDVI) \cite{Rouse} is calculated as a ratio between the sum and the difference of the two Landsat bands - Near-Infrared (NIR) and Red. The NDVI is the most well know and widely used ecological indicator for assessing vegetation conditions due to the optimally created formula that includes the spectral properties of Red and NIR bands. Thus, compared to the other bands, the content of chlorophyll in healthy vegetation is reflected by the NIR band much more. At the same time, it absorbs Red light. Therefore, using these two bands in a well-established combination that gives robust results indicating the presence and amount of chlorophyll in leaves. Such effectiveness of the NDVI resulted in its numerous applications for the environmental analysis \cite{w13192707,rs12244119,GHEBREZGABHER2020249,RULINDA2012169}. The values of the NDVI vary within the range from -1 to +1 where negative values indicate water, lower values above 0 show bare soil and land, and mediocre values (around 1.30-0.50) signify sparsely vegetated areas. Higher values around +1 identify dense vegetation with healthy canopy and green foliage. The NDVI is computed using Equation \ref{eq:1}:

\begin{equation}\label{eq:1}
NDVI=\dfrac{(NIR-Red)}{(NIR+Red)}
\end{equation}

Using R, we computed the NDVI using a script presented in the Appendix, Listing \ref{lst01}. 

\subsubsection{Green Normalized Difference Vegetation Index (GNDVI)}

The GNDVI index is an updated and modified NDVI to remotely sense the presence and vitality of vegetation. Compared to the NDVI, it also uses the NIR Band for computation but replaces the Red Band (i.e., Band 4 for the Landsat 8-9 OLI/TIRS with a spectral band between 640 and 670 nm wavelengths) by the Green Band (Band 3 for the Landsat 8-9 OLI/TIRS with a spectral band between 530 and 590 nm wavelengths) from the visible spectrum. Therefore, the GNDVI estimates the chlorophyll content more precisely compared to the NDVI and enables to analyse photo synthetic activity in plants. The applications of the GNDVI are mostly intended at detecting wilted, ill or aging plants and monitoring plant stress. 

Moreover, it is also suitable for estimating the nitrogen content in the plant leaves, as well as monitoring mature vegetation coverage in the final stages of the crop cycle or vegetation with dense canopy. Similar to the NDVI, the range of values in the GNDVI also varies between -1 and 1, where negative values are associated with the areas of water or bare soil, while higher values indicate healthy vegetation. In this study, the GNDVI was computed due to its technical and practical advantages. Thus, its advantages include the low number of required spectral bands (only NIR and Green, as indicated in the formula of Equation \ref{eq:5}) which results in easy computation in R, and reliability in the analysis of vegetation health due to the more precise estimation of the chlorophyll content compared to the NDVI. The formula for the GNDVI used for computation is given in Equation \ref{eq:5}:

\begin{equation}\label{eq:5}
	GNDVI=\dfrac{(NIR-Green)}{(NIR+Green)}
\end{equation}

For the Landsat 8-9 OLI/TIRS bands, the Green channel corresponds to the Band 3 and Red channel -- to the Band 4. The adjustment so the GNDVI are based on the maximum sensitivity of reflectance of leaves since the reflectance near 670 nm of the Red Band is not sensitive to chlorophyll concentration due to the saturation of the relationship between absorptions and chlorophyll concentration \cite{GITELSON1996289}. Therefore, the replacement of the Red Band by the Green Band results in a more accurate estimate for the concentration of chlorophyll and the rate of photosynthesis. Using R syntax, we computed the GNDVI using script presented in the Appendix, Listing \ref{lst02}. 

\subsubsection{Normalized Difference Water Index (NDWI)}

Developed by Gao in 1996 \cite{GAO1996257} as a complementary to the NDVI, the NDWI is strongly related to the plant water content as it is sensitive to liquid water content of leaves based on the ratio between the difference and sum of NIR and SWIR. The second version of the NDWI developed by \cite{rs5073544} uses a Green band instead of SWIR which can be used to detect surface open water areas as well as evaluate the level of turbidity. In this case, the NDWI is computed using the ratio between the difference and sum of the NIR and visible Green Bands \cite{ZHENG2021129488}. In our study we used the first version developed by Gao following the Equation \ref{eq:2}:

\begin{equation}\label{eq:2}
	NDWI=\dfrac{(NIR-SWIR)}{(NIR+SWIR)}
\end{equation}

The advantages of the NDWI is that it lessens the effects from the reflectance of soil and vegetation foliage which makes it suitable for sandy areas such as in Sahelian Sudan. A special value of the NDWI for environmental monitoring is that it can be used for indicating areas with water deficit \cite{TENG2021107260}. Therefore, the NDWI presents a useful and precise tool in the assessment of the health of vegetation in arid areas prone to droughts, such as Sudanian Sahel, through computing the moisture level in plants as detected water content in leaves. Thus, during the drought periods, foliage is strongly affected by the water deficit which results in plant illness. Specifically for the agricultural areas of Khartoum, it leads to the crop failure and decreases the production which affects the harvest. Detecting the plant water stress is essential to evaluate the current state of plants and make a prognosis for future development, e.g., by indicating the areas that are in need of irrigation. Using the R library 'terra', we computed the NDWI using script presented in the Appendix in Listing \ref{lst03}.

\subsubsection{Optimized Soil Adjusted Vegetation Index (OSAVI)}

Originally developed by \cite{Rondeaux}, the Optimized Soil Adjusted Vegetation Index (OSAVI) is well adjusted to monitor vegetation health specifically in the regions with low vegetation coverage, such as Sahelian region of Sudan due to corrections of bare soil and desert areas. Specifically, the OSAVI uses a defined value of 0.16 as a factor to adjust the canopy background as this value gives higher soil variation compared to the original SAVI. At the same time, the OSAVI has a higher sensitivity to vegetation coverage that covers more than a half of the area in percentage. Due to such characteristics, the OSAVI suitable to monitor the existing productivity changes in agricultural areas near the Khartoum and along the banks of the Nile River. Moreover, the corrections fro the reflectance from bare soil and lands enables the detect areas prone to droughts using comparison of the multi-temporal satellite images. Therefore, the application of the OSAVI is especially suitable in the regions with sparse vegetation, as in the southern Sudan where soil is often visible through the canopy of vegetation. The OSAVI is computed using Equation \ref{eq:4}:

\begin{equation}\label{eq:4}
	OSAVI=\dfrac{(NIR-Red)}{(NIR+Red+0.16)}
\end{equation}

In R language, we computed the OSAVI by script showed in the Appendix, Listing \ref{lst04}.

\subsubsection{Infrared Percentage Vegetation Index (IPVI)}

The values of the Infrared Percentage Vegetation Index (IPVI) depend on the chlorophyll content in leaves. Developed originally by \cite{CRIPPEN199071}, it always has positive values ranging from 0 to +1 and indicates the photosynthetic activity of the canopy cover and total of chlorophyll in leaves. Further, as a modified variation of the NDVI, it is especially suitable to indicate yellowish or shed leaves in the semi-arid areas such as southern Sudan, due to the adjusted spectral regions in the formula. Thus, the chlorophyll content directly depends on the nitrogen level in plants which contributes to the greenness of foliage and is reflected in the optical characteristics of leaves. In such way, the IPVI is suitable for the periods of active vegetation development due to the indication of leaf pigment content and variations showing always positive values above zero, in contrast with the NDVI. Therefore, the IPVI is useful to indicate the yield characterization and stress level in leaves in the regions affected by droughts in the semi-arid ares to the norther of Khartoum. In this way, the IPVI is a robust indicator of the drought severity in southern Sudan, and its effects on the arid and semi-arid ecosystems of Sudanian Sahel. Using the comparison of images covering target test area on the given time periods in 2013, 2018 and 2022, the difference in the vegetation patterns shows the variations in health conditions of leaves. The IPVI is computed using the Equation \ref{eq:3}:

\begin{equation}\label{eq:3}
	IPVI=\dfrac{NIR}{NIR + Red}
\end{equation}

The computation of IPVI using the R library 'terra' was done by the code presented in the Appendix, Listing \ref{lst05}. The application of RStudio version 2023.03.0+386 \cite{RStudio} for image processing included image processing, computing and plotting the vegetation indices index using Landsat 8 OLI/TIRS image. 

%%%%%%%%%%%%%%%%%%%%%%%%%%%%%%%%%%%%%%%%%%
%\newpage
\section{Results}

\hl{All the computations were performed using computer MacBook Air with chip Apple M1 2020, operation system MacOS Ventura version 13.2.1 which uses 8GB of superfast unified memory. On this machine, the technical advantages and capabilities of R and its geo-processing packages 'terra' and 'raster' were evaluated in terms of processing time. We quantified several computational aspects, which demonstrated the following results. 1. Time taken to read data: less then 1 second (i.e. the data were read in instantly as soon as the 'enter' was pressed); 2. Time taken to process data (seconds): 9 seconds to process the following commands corresponding to the lines 7 to 20 in the case of Listing 1 R code for computing the NDVI): vi <- function(); filenames <- paste0(); landsat <- rast(); ndvi <- vi(landsat, 5, 4); options(scipen=10000); and colors <- brewer.rdylgn(100). Next, the graphical plotting of the image (lines 21-22 in example of the same Listing) was perform during 2 seconds, and plotting the histogram (lines 24-27) took 1 second. 4. Overall processing speed took seconds): 12 seconds for each processed Landsat OLI/TIRS image. 3. Memory utilization (Mb) demonstrated effective parameters due to the Mac with chip Apple M is powerful in terms of RAM. Therefore, the highest memory consumption was 60 Mb for 15 maps which includes the computer 5 vegetation indices for 3 different calendar years.}

We provide in Table \ref{tab02} for each set of images their VI values for years 2013, 2018 and 2022, derived from the R report on data analysis, as it appears in the output console. These values were collected and summarised in Table \ref{tab02} for comparison of variations in the VI's values. In the second to fourth columns we indicate the years of observations and output results obtained by R 'terra' library algorithm.  For each image, we collected the data's minimal and maximal values to assess its global range. Thus, the results of the VI values derived from the Landsat 8 OI/TIRS images over Khartoum region area are presented in Table \ref{tab02}. These data show the variations indicating changes in vegetation health. The maximal values indicate the healthy green vegetation which corresponds to greenness due to the chlorophyll content in leaves. The changes between the 2013 and 2018 show the decline for the values in the following indices: NDVI, GNDVI, NDWI, OSAVI and a vey slight increase in the IPVI. The analysis of data on 2018 and 2022 shows that there is a slight rise in values for the same indices expect for the IPVI which demonstrates a slight decrease.

\begin{table}[H] 
\caption{Evaluation results of the VI values derived from the Landsat 8 OLI/TIRS images over the Khartoum region area (2013, 2018 and 2022).\label{tab02}}
\newcolumntype{C}{>{\centering\arraybackslash}X}
\begin{tabularx}{\textwidth}{CCCC}
\toprule
\textbf{VI}	& \textbf{2013}	& \textbf{2018} & \textbf{2022}\\
\midrule
	$NDVI_{min}$ & -0.2742168 & -0.2817099 & -0.2791084 \\
	$NDVI_{max}$ & 0.7246868 & 0.5796811 & 0.6044010 \\\hline
	$GNDVI_{min}$ & -0.2387020 & -0.2830099 & -0.2253696 \\
	$GNDVI_{max}$ & 0.6371941 & 0.5690726 & 0.7235179 \\\hline
	$NDWI_{min}$ & -0.4046338 & -0.5656914 & -0.5743572 \\
	$NDWI_{max}$ & 0.3880472 & 0.3764042 & 0.4152763\\\hline
	$OSAVI_{min}$ & -0.2742153 & -0.2817084 & -0.2791067 \\
	$OSAVI_{max}$ & 0.7246833 & 0.5796785 & 0.6043976 \\\hline
	$IPVI_{min}$ & 0.1376566 & 0.2101595 & 0.1977995 \\
	$IPVI_{max}$ & 0.6371084 & 0.6408550 & 0.6395542 \\
\bottomrule
\end{tabularx}
\end{table}

\subsection{Normalized Difference Vegetation Index (NDVI)}
The images of the NDVI (Figure \ref{fig04}) allow the differentiation of vegetation areas or cultivated areas from soil in the region southwards from the Khartoum and along the flow of the Greater Nile. The negative NDVI values were classified as non-vegetation (i.e. soil and bare land), and positive NDVI values -- as vegetation. 

\begin{figure}[H]
\makebox[0.99\linewidth][r]{%
	\begin{subfigure}[b]{.55\textwidth}
		\centering
			\hspace{30pt}\includegraphics[width=0.95\textwidth]{Fig_04a.jpg}
		\caption{}
	\end{subfigure}\hspace*{\fill}%
	\begin{subfigure}[b]{.55\textwidth}
		\centering
			\hspace{30pt}\includegraphics[width=0.95\textwidth]{Fig_04b}
		\caption{}
	\end{subfigure}%
}\\
\vfill \vspace{1mm}
\makebox[0.99\linewidth][r]{%
	\begin{subfigure}[b]{.55\textwidth}
		\centering
			\hspace{30pt}\includegraphics[width=0.95\textwidth]{Fig_04c.jpg}
		\caption{}
	\end{subfigure}%
	\begin{subfigure}[b]{.55\textwidth}
		\centering
			\hspace{30pt}\includegraphics[width=0.95\textwidth]{Fig_04d}
		\caption{}
	\end{subfigure}%
}\\
\vfill \vspace{1mm}
\makebox[0.99\linewidth][r]{%
	\begin{subfigure}[b]{.55\textwidth}
		\centering
			\hspace{30pt}\includegraphics[width=0.95\textwidth]{Fig_04e.jpg}
		\caption{}
	\end{subfigure}%
	\begin{subfigure}[b]{.55\textwidth}
		\centering
			\hspace{30pt}\includegraphics[width=0.95\textwidth]{Fig_04f}
		\caption{}
	\end{subfigure}%
}
\caption{NDVI computed from the Landsat 8 OLI/TIRS images on confluence of the White Nile and Blue Nile, Sudan for three years (always December): (\textbf{a}) 2013, (\textbf{b}) 2018, (\textbf{c}) 2022, (\textbf{d}) histogram for 2013, (\textbf{e}) histogram for 2018, (\textbf{f}) histogram for 2022.}\label{fig04}
\end{figure}

The highest NDVI values (0.80 to 1.0) identify forest areas. For the desert areas with dominating sands such as Sudan, the NDVI values do not exceed 0.50. The minimal values determined according to the histograms obtained from the NDVI images for years 2013, 2018 and 2022 were 0.27, 0.28 and 0.28 (Table \ref{tab02}), whereas the highest determined values were 0.72, 0.58 and 0.60 for the same years. The minimum and maximum values were derived from the computed vegetation indices as a statistical summary of data analysis. They indicate the range in greenness in the vegetation: the higher values correspond to the healthy green vegetation while lower values show bare land. The decrease in higher values shows the decline in green and healthy plants. 

The pixels that were not classified to either the soil, water or vegetation classes, were identified as sand desert areas (mostly beige-coloured in the images in Figure \ref{fig04}) pixels due to low NDVI values around 0.2 or less. The lowest NDVI values less then 0.1 are classified as rocky surface and bare land. The identification method of soil/water and vegetation is based on the characteristics of the NDVI which ranges from -1.0 to 1.0 with negative values indicating clouds and water, positive values near zero ($<0.1$) -- bare soil, barren rock or sand in deserts. Shrubs, grasslands and other types of sparse vegetation have moderate NDVI values between 0.2 to 0.5. Higher positive values of NDVI correspond to dense green vegetation ($>0.6 and <0.9$) which correspond to tropical forests or crops at their highest growth stage. Light ivory colour corresponds to the patterns of the dry creeks or beds of wadi ephemeral streams. The minimal and maximal values of the NDVI varied for years 2013, 2018 and 2022 are as follows: from -0.2742168 to 0.7246868 in 2013, from -0.2817099 to 0.5796811 in 2018, and from -0.2791084 to 0.6044010 in 2022. As one can note, the maximal values decreased from 2013, which indicates the general trend in desertification, however, from 2018 to 2022, there is a slight increase again in the highest values which can be explained by increased precipitation and fluctuations in Nile hydrology. 

\subsection{Green Normalized Difference Vegetation Index (GNDVI)}

The highest values of the GNDVI did not exceed 0.72 for the area of Sudan due to the scarce vegetation coverage and dominating desert areas, Figure \ref{fig05}. The negative values are associated with the areas of water or bare soil, while higher values indicate healthy vegetation and selected agricultural croplands along the banks of the White, Blue and Greater Niles. As seen from Figure \ref{fig05}, the GNDVI captured plant greenness with changes in vegetation health and greenness corresponding to the colour on the images varying from bright green to yellow on sequence of images for 2013, 2018 and 2022. In this case, the variations in shadows of colours indicate the differences in plant vigour which is a measure of the increase in plant growth or foliage volume through time after planting \cite{Dictionary}. Moreover, it indicate at the plant health which correlates to the irrigated lands with higher level of moisture in soil that can be distinguished against drylands and desert areas. 

Further comparing the values in Table \ref{tab02}, the range of values for the GNDVI is narrower than that of the NDVI: with minimal -0.24, -0.28 and -0.23 for years 2013, 2018 and 2022, while the highest values are 0.64, 0.57 and 0.72 for the same years. Accordingly, with greener plants visualised in the GNDVI, Figure \ref{fig05} shows lower and better distinguished values. Moreover, the histograms present more diversified range of values for various land classes: water, bare soil and sand, rocky areas of deserts and vegetation with agricultural areas and natural aquatic plants in the Nile rivers. Finally, the slight increase in the GNDVI values indicating photo synthetic activity in leaves is affected by changes in water and N absorption in plant that includes aquatic vegetation in the bank areas of the White Nile. 

\begin{figure}[H]
\makebox[0.99\linewidth][r]{%
	\begin{subfigure}[b]{.55\textwidth}
		\centering
			\includegraphics[width=0.95\textwidth]{Fig_05a}
		\caption{2013}
	\end{subfigure}%
	\begin{subfigure}[b]{.55\textwidth}
		\centering
			\includegraphics[width=0.95\textwidth]{Fig_05b}
		\caption{2013: histogram}
	\end{subfigure}%
}\\
\vfill \vspace{1mm}
\makebox[0.99\linewidth][r]{%
	\begin{subfigure}[b]{.55\textwidth}
		\centering
			\includegraphics[width=0.95\textwidth]{Fig_05c}
		\caption{2018}
	\end{subfigure}%
	\begin{subfigure}[b]{.55\textwidth}
		\centering
			\includegraphics[width=0.95\textwidth]{Fig_05d}
		\caption{2018: histogram}
	\end{subfigure}%
}\\
\vfill \vspace{1mm}
\makebox[0.99\linewidth][r]{%
	\begin{subfigure}[b]{.55\textwidth}
		\centering
			\includegraphics[width=0.95\textwidth]{Fig_05e}
		\caption{2022}
	\end{subfigure}%
	\begin{subfigure}[b]{.55\textwidth}
		\centering
			\includegraphics[width=0.95\textwidth]{Fig_05f}
		\caption{2022: histogram}
	\end{subfigure}%
}
\caption{GNDVI computed from the Landsat 8 OLI/TIRS images on confluence of the White Nile and Blue Nile, Sudan for three years (always December): (\textbf{a}) 2013, (\textbf{b}) 2018, (\textbf{c}) 2022, (\textbf{d}) histogram for 2013, (\textbf{e}) histogram for 2018, (\textbf{f}) histogram for 2022.}\label{fig05}
\end{figure}

\subsection{Normalized Difference Water Index (NDWI)}

The interpretation of the values on the resulting images based on the NDWI (Figure \ref{fig06}) are as follows: the negative values (from -1 up to 0) signify areas with no vegetation or water content, while those above 0 show surfaces with water content where the higher moisture level correlates with higher values of the NDWI, respectively. The NDWI values vary between -1 to +1, which depends on the total amount of water in leaves, type of vegetation and percentage of the foliage coverage. Thus, sparse vegetation in semi-arid regions of the Sudanian Sahel will naturally have a lower value of the NDWI compared to the regions of dominating forests with dense canopy. The analysis of the NDWI images reveals contours of the dry creeks in the desert areas with stream patterns and beds of drainage courses that contribute to the differences in moisture content in soil (Figure \ref{fig06}). 

\begin{figure}[H]
\makebox[0.99\linewidth][r]{%
	\begin{subfigure}[b]{.5\textwidth}
		\centering
			\includegraphics[width=0.95\textwidth]{Fig_06a}
		\caption{2013}
	\end{subfigure}%
	\begin{subfigure}[b]{.5\textwidth}
		\centering
			\includegraphics[width=0.95\textwidth]{Fig_06b}
		\caption{2013: histogram}
	\end{subfigure}%
}\\
\vfill \vspace{1mm}
\makebox[0.99\linewidth][r]{%
	\begin{subfigure}[b]{.5\textwidth}
		\centering
			\includegraphics[width=0.95\textwidth]{Fig_06c}
		\caption{2018}
	\end{subfigure}%
	\begin{subfigure}[b]{.5\textwidth}
		\centering
			\includegraphics[width=0.95\textwidth]{Fig_06d}
		\caption{2018: histogram}
	\end{subfigure}%
}\\
\vfill \vspace{1mm}
\makebox[0.99\linewidth][r]{%
	\begin{subfigure}[b]{.5\textwidth}
		\centering
			\includegraphics[width=0.95\textwidth]{Fig_06e}
		\caption{2022}
	\end{subfigure}%
	\begin{subfigure}[b]{.5\textwidth}
		\centering
			\includegraphics[width=0.95\textwidth]{Fig_06f}
		\caption{2022: histogram}
	\end{subfigure}%
}
\caption{NDWI computed from the Landsat 8 OLI/TIRS images on confluence of the White Nile and Blue Nile, Sudan for three years (always December): (\textbf{a}) 2013, (\textbf{b}) 2018, (\textbf{c}) 2022, (\textbf{d}) histogram for 2013, (\textbf{e}) histogram for 2018, (\textbf{f}) histogram for 2022.}\label{fig06}
\end{figure}

Such patterns are indicative of moist wadi sediments and associated vegetation strips. The range of values in the NDWI shows slow dynamics in the lowest values that is gradually increasing from -0.40 in 2013, to -0.56 in 2018 and 0.57 in 2022, along with the raising highest values: 0.38 in 2013, 0.38 in 2018 and 0.42 in 2022. The values of the NDWI rise along with the increased leaf water content as well as vegetation fraction cover. Therefore, higher correspond to high water content in leaves and to increase in canopy. The causes for the increase in water content in leaves include climate and hydrological factors such as higher precipitation level and increased water discharge from the Nile. Since the NDWI is specifically designed to monitor vegetation in drought affected areas such as Sudan, such variations indicate at the increase in both the content of vegetation water because of the strong chlorophyll absorption in SWIR and the structure in plant canopy. 

On the one hand, partial vegetation coverage in the desert areas of central Sudan results in soil effects on the NDWI values, especially for the landscapes with interspersed soils and vegetation patterns. Thus, when the leaf layer in the canopy decreases, the absolute values of reflectances in the spectral region of the NDWI are decreased which results in lower values. The range of the NDWI values in Khartoum is narrow which is seen on the histograms with most of the values varying from -0.20 to +0.20. More specifically, the interpretation of the values of the NDWI follows the ranges: positive values from +0,2 to +1.0 -- water area; 0.0 to +0,2 -- flooded areas and high humidity, -0,3 to 0.0: moderate drought soil and non-water surfaces, -1.0 to -0.3: drylands, drought and non-water surfaces. The NDWI indicates the moisture level of vegetation which is suitable for the sandy areas in the Sahelian Sudan.

%--------------------->
\subsection{Optimized Soil Adjusted Vegetation Index (OSAVI)}

The computed Optimized Soil Adjusted Vegetation Index (OSAVI) is illustrated in Figure \ref{fig07}). 

\begin{figure}[H]
\makebox[0.99\linewidth][r]{%
	\begin{subfigure}[b]{.5\textwidth}
		\centering
			\includegraphics[width=0.95\textwidth]{Fig_07a}
		\caption{2013}
	\end{subfigure}%
	\begin{subfigure}[b]{.5\textwidth}
		\centering
			\includegraphics[width=0.95\textwidth]{Fig_07b}
		\caption{2013: histogram}
	\end{subfigure}%
}\\
\vfill \vspace{1mm}
\makebox[0.99\linewidth][r]{%
	\begin{subfigure}[b]{.5\textwidth}
		\centering
			\includegraphics[width=0.95\textwidth]{Fig_07c}
		\caption{2018}
	\end{subfigure}%
	\begin{subfigure}[b]{.5\textwidth}
		\centering
			\includegraphics[width=0.95\textwidth]{Fig_07d}
		\caption{2018: histogram}
	\end{subfigure}%
}\\
\vfill \vspace{1mm}
\makebox[0.99\linewidth][r]{%
	\begin{subfigure}[b]{.5\textwidth}
		\centering
			\includegraphics[width=0.95\textwidth]{Fig_07e}
		\caption{2022}
	\end{subfigure}%
	\begin{subfigure}[b]{.5\textwidth}
		\centering
			\includegraphics[width=0.95\textwidth]{Fig_07f}
		\caption{2022: histogram}
	\end{subfigure}%
}
\caption{OSAVI computed from the Landsat 8 OLI/TIRS images on confluence of the White Nile and Blue Nile, Sudan for three years (always December): (\textbf{a}) 2013, (\textbf{b}) 2018, (\textbf{c}) 2022, (\textbf{d}) histogram for 2013, (\textbf{e}) histogram for 2018, (\textbf{f}) histogram for 2022.}\label{fig07}
\end{figure}

The OSAVI was the highest for 2013 with the maximal reached value of 0.72 (Figure \ref{fig07}), after which it dropped to 0.58 in 2018 indicating the decease in the leaf chlorophyll content reflected in the OSAVI values; since 2018 to 2022 the values stabilised and slightly increased to 0.60 due to the effects from the White Nile hydrology and climatic variations. Being a modified index derived from the NDVI, the OSAVI also has a range of values varying between the -0.28 to +0.72, with negative values indicating water areas of the White and Blue Niles, lower areas slightly above 0 and lower than 0.2 indicating bare land, sandy areas, soil with no or very sparse vegetation, while higher values close to +0.50 showing healthy vegetation and higher chlorophyll content in leaves. The OSAVI values typically range from -1 to +1 where high values indicate denser, healthier vegetation, lower values indicate less vigour and medium values corresponding to the agricultural lands. The highest and minimum values in this region show the range in data which indicate the changes in the greenness of vegetation and health of plant leaves.

The OSAVI image shows that the land classes with values between 0.30 and 0.40 representing sparse vegetation, 0.40 to 0.50 representing medium vegetated area typically for agriculture lands and those above 0.50 representing occasional regions with high vegetated area around the Blue Nile. The analysis of the OSAVI and comparison with other indices reveals that the OSAVI classifies pixels on the image was more detailed compared to the NDVI and GNDVI as related to resolution. Thus, it identifies vegetation better using spectral reflectance within each cluster of data range which results in the detection of the biomass and leaf foliage areas in the regions covered by sandy soils and semi-arid lands such as Sudan. 
%--------------------->
\subsection{Infrared Percentage Vegetation Index (IPVI)}

The majority of values of the Infrared Percentage Vegetation Index (IPVI) for the area of Sudan, Figure \ref{fig08}, lies within a very narrow interval of 0.25 to 0.48 which matches the dominance of sandy dune areas and bare soils in the surrounding deserts. Note that all the values of IPVI are never negative and always positive due to the formula used in computation, which distinguishes this index from the previous ones.

\begin{figure}[H]
\makebox[0.99\linewidth][r]{%
	\begin{subfigure}[b]{.55\textwidth}
		\centering
			\includegraphics[width=0.95\textwidth]{Fig_08a}
		\caption{2013}
	\end{subfigure}%
	\begin{subfigure}[b]{.55\textwidth}
		\centering
			\includegraphics[width=0.95\textwidth]{Fig_08b}
		\caption{2013: histogram}
	\end{subfigure}%
}\\
\vfill \vspace{1mm}
\makebox[0.99\linewidth][r]{%
	\begin{subfigure}[b]{.55\textwidth}
		\centering
			\includegraphics[width=0.95\textwidth]{Fig_08c}
		\caption{2018}
	\end{subfigure}%
	\begin{subfigure}[b]{.55\textwidth}
		\centering
			\includegraphics[width=0.95\textwidth]{Fig_08d}
		\caption{2018: histogram}
	\end{subfigure}%
}\\
\vfill \vspace{1mm}
\makebox[0.99\linewidth][r]{%
	\begin{subfigure}[b]{.55\textwidth}
		\centering
			\includegraphics[width=0.95\textwidth]{Fig_08e}
		\caption{2022}
	\end{subfigure}%
	\begin{subfigure}[b]{.55\textwidth}
		\centering
			\includegraphics[width=0.95\textwidth]{Fig_08f}
		\caption{2022: histogram}
	\end{subfigure}%
}
\caption{IPVI computed from the Landsat 8 OLI/TIRS images on confluence of the White Nile and Blue Nile, Sudan for three years (always December): (\textbf{a}) 2013, (\textbf{b}) 2018, (\textbf{c}) 2022, (\textbf{d}) histogram for 2013, (\textbf{e}) histogram for 2018, (\textbf{f}) histogram for 2022.}\label{fig08}
\end{figure}

%%%%%%%%%%%%%%%%%%%%%%%%%%%%%%%%%%%%%%%%%%
\section{Discussion}

The method of R \hl{'raster' and 'terra' libraries} is advantageous for \hl{satellite image processing aimed at computing vegetation indices. The computations were conduced on a MacBook Air PC (i.e., chip Apple M1 2020, MacOS Ventura v. 13.2.1, 8GB of main memory). The evaluated processing speed reported in the Results section of this manuscript demonstrated technical advantages and benefits of R and its libraries 'terra' and 'raster'. The performance of R libraries was evaluated in terms of processing time with quantified computational aspects. High speed of computations demonstrated by R is beneficial for }repetitive workflow of image processing when several scenes need to be processed sequentially for several years. \hl{Specifically, high speed of reading the data and total time of image processing (12 seconds for each image) confirmed that R is advantageous for remote sensing data processing, computational and mapping tasks.}

\hl{Moreover, using} scripts\hl{ }enables to save time through automation of the repetitive tasks\hl{. This is more advantageous compared to the GIS with traditional GUI interface used for satellite image processing, such as ENVI GIS, SAGA GIS or Erdas IMAGINE}. This is especially useful \hl{for a} computation of several\hl{ }images to calculate several VIs using \hl{embedded }formulae\hl{. }Further, R scripts eliminate errors when processing the satellite images due to high level of automation. Applying scripts \hl{customized} for each scene \hl{increases the speed of data processing which is advantageous when working with several images.} For instance, scripts call embedded functions to perform parts of computations and data processing. Moreover, scripts provide quick plotting of maps and graphics that visualise the results. This help to highlight the changes in vegetation patterns when maps visualise the VIs on several years for comparison. 

\hl{A }package-based approach of R that \hl{is} used \hl{in this paper }as \hl{a }target tool to \hl{remote sensing data} processing \hl{which enables to perform computation of VI for analysis of satellite images. The presented use of R libraries shows the usefulness of this approach as an analytical tool with which we computed several VI and thus revealed changes in vegetation patterns of southern Sudan since 2013}. \hl{The use of R presents }opportunities for sustainable environmental monitoring of Sudan through integration of the remotely sensed data and advanced programming technologies. In this paper, we have proposed the application of R language for computation of several vegetation indices using the 'terra' library and satellite images Landsat 8-9 OLI/ITIRS. Compared to traditional GIS approaches, the proposed method demonstrates high efficiency and effectiveness for geospatial data processing aimed at environmental analysis. In the performed experiments in calculating VIs over the Landsat 8 OLI/TIRS data, our study has shown that, in most cases of the indices, a general trend in decrease of VI values is notable for data on 2013, 2018 and 2022, as shown in Table \ref{tab02} and visualised in Figures \ref{fig04}, \ref{fig05}, \ref{fig06}, \ref{fig07} and \ref{fig08}. This is clearly visible from the results obtained using the estimated parameters on indices NDVI, OSAVI, IPVI. 

Using three satellite images Landsat 8-9 OLI/ITIRS collected for the region of Khartoum area, we detected a slight decrease in vegetation greenness from 2013 to 2022 based on the comparison of the dynamics of the satellite-derived VIs. For each image, the computation of five of the presented VIs was performed for three corresponding years by applying our R scripts. The visualisation of these indices demonstrates a function of vegetation health, greenness and vigour in the conditions of the Sahelian climate of Sudan with variations over the analysed years. The proposed R implementation shows promising results when detecting anomalies in vegetation conditions in southern Sudan, as the VIs provide essential information on plant greenness reflected as chlorophyll content. The chlorophyll content relates to the water content in leaves which in turn is controlled by regional climatic setting and hydrological processes of the Nile River and its tributaries - the White and Blue Niles. 

We have addressed the processing of satellite images by computing a series of vegetation indices for the regions prone to desertification in the Sahelian area of Sudan. Our approach represents each Landsat 8-9 OLI/TIRS image with five derived vegetation indices, each tuned for various environmental purposes, and compared them for the years 2013, 2018 and 2022: detecting dense vegetation with healthy canopy and green foliage (NDVI), estimation of the canopy background in the regions with sparse vegetation and low percentage of foliage coverage as typical for Sudan (OSAVI), evaluating the maturity of the canopy based on the leaf pigment content and nitrogen level in plants (IPVI), evaluation of concentration of chlorophyll in leaves and the rate of photosynthesis (GNDVI), indicating areas with water deficit (NDWI). In this way, the comparative analysis of indices visualised for various periods can be used to identify vegetation conditions aimed at the environmental monitoring of the Sahel region of Northeast Africa.

We presented further insights into raster image processing by R scripting and optimization methods for a time series analysis of several scenes. The computational workflow has demonstrated that the proposed algorithms of R perform is far better in terms of the processing and analysis of spatial datasets. The experiments have demonstrated that scientific scripting in the R language plays an important role in geospatial image processing, visualisation and analysis, since it enables integrating various components or libraries including graphical packages used for mapping and plotting as well as packages for spatial data processing which can be integrated into a coherent scripting workflow.  

%%%%%%%%%%%%%%%%%%%%%%%%%%%%%%%%%%%%%%%%%%
\section{Conclusions}

In this paper, we highlighted that the performance of the R packages is effective for a framework of satellite image processing and enables to compute the vegetation indices. With the presented maps of the VIs made based on the Landsat OLI/TIRS images, we observe the changes in the vegetation patterns in the test area. We further showed that there is a trend to the growth and greenness of vegetation over the period of 2013 to 2022 in the study area in the confluence of the White and Blue Niles in southern Sudan, which is caused by the rising temperature and decreased precipitation due to the climatic and environmental impacts effects of desertification.

The Landsat 8-9 OLI/TIRS images contain a complex image texture with distinct values of pixels due to different spectral reflectance of vegetation and other land cover types. Five VIs of vegetation distribution analysis have been applied and brought to practical applicability by means of R language. Specifically, we tested five VIs for the Landsat 8-9 OLI/TIRS images and compared the results for the following indices: 1) Normalized Difference Vegetation Index (NDVI), 2) Normalized Difference Water Index (NDWI), 3) Infrared Percentage Vegetation Index (IPVI), 4) Optimized Soil Adjusted Vegetation Index (OSAVI) and 5) Green Normalized Difference Vegetation Index (GNDVI). 

The computational workflow of R libraries included extraction of image sensor DN values of pixels, computing the data range with min/max values, classifying the raster object and visualising the image using selected colour schemes. Within all these methods, data processing velocity can be estimated as higher compared to conventional GIS approaches, which offers additional possibilities of programming for geospatial data analysis and image processing. As such, we demonstrated the usefulness of the R programming language for machine-based information extraction from satellite images. We also showed the effectiveness of the remote sensing data as a source of information derived from Earth observation satellites for regional studies focused on Northeastern Africa.

Image analysis in Earth sciences has made big progress due to the introduction of programming applications and scripts that significantly automate the routine of data processing. In the near future we can expect the possibility of stream processing applied for large collections of satellite images as a series of data for operative monitoring. Thus, further environmental applications of image processing could for instance include, the monitoring of burned areas in savannah fires using satellite scenes processed straight after they are generated by sensors. Furthermore, the use of scripts for image processing opens the way for the geospatial analysis of diverse objects and structures of very high complexity that constitute the landscapes of the Earth. This includes the detailed diversification of land cover types, vegetation associations and plant communities that can be identified from images.

In monitoring the Earth's landscapes, it is important to not only use topographic data, but also the environmental characteristics that reflect the complex interactions from climate and hydrological parameters, as well as vegetation responses to the external climatic effects. These can all be derived from the satellite images such as Landsat 8-9 OLI/TIRS using advanced methods of image processing as demonstrated in this paper. Future research will be focused on developing further approaches using scripting algorithms that deal with satellite images to detect other parameters of land surface and vegetation coverage for environmental mapping, such as brightness and saturation indices, or evaluating other vegetation indices for detecting the salinity of soil in arid lands of Africa. Such characteristics reflect the climatic-vegetation interactions that can be used for operational environmental monitoring based on the spaceborne data.  

%%%%%%%%%%%%%%%%%%%%%%%%%%%%%%%%%%%%%%%%%%
\authorcontributions{Supervision, conceptualization, methodology, software, resources, funding acquisition, and project administration, O.D.; writing---original draft preparation, methodology, software, data curation, visualization, formal analysis, validation, writing---review and editing, and investigation, P.L. All authors have read and agreed to the published version of the manuscript.}

\funding{The publication was funded by the Editorial Office of Journal of Imaging, Multidisciplinary Digital Publishing Institute (MDPI), by providing 100\% discount for the APC of this manuscript. This project was supported by the Federal Public Planning Service Science Policy or Belgian Science Policy Office, Federal Science Policy - BELSPO (B2/202/P2/SEISMOSTORM).}

\institutionalreview{Not applicable.}

\informedconsent{Not applicable.}

\dataavailability{Not applicable} 

\acknowledgments{The authors thank the reviewers for reading and review of this manuscript.}

\conflictsofinterest{The authors declare no conflict of interest.} 

\abbreviations{Abbreviations}{
The following abbreviations are used in this manuscript:\\

\noindent 
\begin{tabular}{@{}ll}
GMT & Generic Mapping Tools \\
GNDVI & Green Normalized Difference Vegetation Index \\
IPVI & Infrared Percentage Vegetation Index \\
Landsat 8-9 OLI/TIRS & Landsat 8-9 Operational Land Imager (OLI) and Thermal Infrared (TIRS) \\
NDVI & Normalized Difference Vegetation Index \\
NDWI & Normalized Difference Water Index \\
OSAVI & Optimized Soil Adjusted Vegetation Index \\
SAVI & Soil Adjusted Vegetation Index \\
USGS & United States Geological Survey \\
UTM & Universal Transverse Mercator \\
VIs & Vegetation Indices \\
WGS84 & World Geodetic System 1984
\end{tabular}
}

%% Optional
\appendixtitles{no} % Leave argument "no" if all appendix headings stay EMPTY (then no dot is printed after "Appendix A"). If the appendix sections contain a heading then change the argument to "yes".
\appendixstart
\appendix
\section[\appendixname~\thesection]{Metadata on the Landsat 8-9 OLI/TIRS images (USGS EarthExplorer).}

\begin{table}[H]
\scriptsize
\caption{Metadata on the Landsat OLI/TIRS satellite images used in this study \textsuperscript{1}.\label{tabApp}}
\begin{adjustwidth}{-\extralength}{0cm}
\begin{tabularx}{\fulllength}{l|l|l|l}
    \hline
        \textbf{Data Set Attribute} & \textbf{Attributes for image on 2022/12/21} & \textbf{Attributes for 2018/12/18} & \textbf{Attributes for 2013/12/20} \\ \hline
        Landsat Scene Identifier & LC91730492022355LGN01 & LC81730492018352LGN00 & LC81730492013354LGN01 \\ \hline
        Date Acquired & 2022/12/21 & 2018/12/18 & 2013/12/20 \\ \hline
        Collection Category & T1 & T1 & T1 \\ \hline
        Collection Number & 2 & 2 & 2 \\ \hline
        WRS Path & 173 & 173 & 173 \\ \hline
        WRS Row & 49 & 49 & 49 \\ \hline
        Target WRS Path & 173 & 173 & 173 \\ \hline
        Target WRS Row & 49 & 49 & 49 \\ \hline
        Nadir/Off Nadir & NADIR & NADIR & NADIR \\ \hline
        Roll Angle & -0.001 & 0.000 & 0.000 \\ \hline
        Date Product Generated L2 & 2023/03/17 & 2020/08/30 & 2020/09/12 \\ \hline
        Date Product Generated L1 & 2023/03/17 & 2020/08/30 & 2020/09/12 \\ \hline
        Start Time & 2022-12-21 08:09:26 & 2018-12-18 08:08:54.736348 & 2013-12-20 08:10:29.988456 \\ \hline
        Stop Time & 2022-12-21 08:09:58 & 2018-12-18 08:09:26.506347 & 2013-12-20 08:11:01.758452 \\ \hline
        Station Identifier & LGN & LGN & LGN \\ \hline
        Day/Night Indicator & DAY & DAY & DAY \\ \hline
        Land Cloud Cover & 0.00 & 0.00 & 0.00 \\ \hline
        Scene Cloud Cover L1 & 0.00 & 0.00 & 0.00 \\ \hline
        Ground Control Points Model & 511 & 521 & 639 \\ \hline
        Ground Control Points Version & 5 & 5 & 5 \\ \hline
        Geometric RMSE Model & 7.263 & 6.387 & 4.378 \\ \hline
        Geometric RMSE Model X & 5.055 & 4.612 & 2.951 \\ \hline
        Geometric RMSE Model Y & 5.216 & 4.418 & 3.234 \\ \hline
        Processing Software Version & LPGS\_16.2.0 & LPGS\_15.3.1c & LPGS\_15.3.1c \\ \hline
        Sun Elevation L0RA & 44.22152569 & 44.38827439 & 44.40055924 \\ \hline
        Sun Azimuth L0RA & 148.68774152 & 148.91482828 & 149.05309430 \\ \hline
        TIRS SSM Model & N/A & FINAL & ACTUAL \\ \hline
        Data Type L2 & OLI\_TIRS\_L2SP & OLI\_TIRS\_L2SP & OLI\_TIRS\_L2SP \\ \hline
        Sensor Identifier & OLI\_TIRS & OLI\_TIRS & OLI\_TIRS \\ \hline
        Satellite & 9 & 8 & 8 \\ \hline
        Product Map Projection L1 & UTM & UTM & UTM \\ \hline
        UTM Zone & 36 & 36 & 36 \\ \hline
        Datum & WGS84 & WGS84 & WGS84 \\ \hline
        Ellipsoid & WGS84 & WGS84 & WGS84 \\ \hline
        Scene Center Lat DMS & "15°54'03.06""N" & "15°54'03.56""N" & "15°54'03.42""N" \\ \hline
        Scene Center Long DMS & "33°06'32.69""E" & "33°07'35.51""E" & "33°05'59.28""E" \\ \hline
        Corner Upper Left Lat DMS & "16°56'34.84""N" & "16°56'44.88""N" & "16°56'44.41""N" \\ \hline
        Corner Upper Left Long DMS & "32°02'24.68""E" & "32°03'25.49""E" & "32°01'44.08""E" \\ \hline
        Corner Upper Right Lat DMS & "16°56'30.70""N" & "16°56'40.13""N" & "16°56'40.63""N" \\ \hline
        Corner Upper Right Long DMS & "34°10'42.96""E" & "34°11'43.87""E" & "34°10'12.61""E" \\ \hline
        Corner Lower Left Lat DMS & "14°50'38.76""N" & "14°50'39.01""N" & "14°50'38.58""N" \\ \hline
        Corner Lower Left Long DMS & "32°03'00.32""E" & "32°04'00.55""E" & "32°02'20.18""E" \\ \hline
        Corner Lower Right Lat DMS & "14°50'35.16""N" & "14°50'34.87""N" & "14°50'35.34""N" \\ \hline
        Corner Lower Right Long DMS & "34°09'59.18""E" & "34°10'59.41""E" & "34°09'29.09""E" \\ \hline
        Scene Center Latitude & 15.90085 & 15.90099 & 15.90095 \\ \hline
        Scene Center Longitude & 33.10908 & 33.12653 & 33.09980 \\ \hline
        Corner Upper Left Latitude & 16.94301 & 16.94580 & 16.94567 \\ \hline
        Corner Upper Left Longitude & 32.04019 & 32.05708 & 32.02891 \\ \hline
        Corner Upper Right Latitude & 16.94186 & 16.94448 & 16.94462 \\ \hline
        Corner Upper Right Longitude & 34.17860 & 34.19552 & 34.17017 \\ \hline
        Corner Lower Left Latitude & 14.84410 & 14.84417 & 14.84405 \\ \hline
        Corner Lower Left Longitude & 32.05009 & 32.06682 & 32.03894 \\ \hline
        Corner Lower Right Latitude & 14.84310 & 14.84302 & 14.84315 \\ \hline
        Corner Lower Right Longitude & 34.16644 & 34.18317 & 34.15808 \\ \hline
%    \end{tabular}
 \end{tabularx}
 \end{adjustwidth}
\end{table}

\section[\appendixname~\thesection]{R scripts used for computing and plotting the vegetation indices}
\subsection[\appendixname~\thesubsection]{R code for computing the NDVI}

\begin{lstlisting}[language=R,caption=R code for computing the NDVI (here: a case of 2022),style=mystyle,label={lst01}]
library(terra)
library(RColorBrewer)
library(Hmisc)
library(pals)
setwd("/Users/polinalemenkova/Documents/R/52_Image_Processing/Sudan_NDVI")
# 1. Normalized Difference Vegetation Index (NDVI) = (NIR - R) / (NIR + R), i.e., NDVI = (Band 5 - Band 4) / (Band 5 + Band 4).
vi <- function(img, k, i) {
  bk <- img[[k]]
  bi <- img[[i]]
  vi <- (bk - bi) / (bk + bi)
  return(vi)
}
# For Landsat NIR = 5, red = 4.
filenames <- paste0('LC09_L2SP_173049_20221221_20221224_02_T1_SR_B', 1:7, ".tif")
filenames
landsat <- rast(filenames)
landsat
ndvi <- vi(landsat, 5, 4)
options(scipen=10000)
colors <- brewer.rdylgn(100)
plot(ndvi, col=colors, font.main = 1, main = "NDVI for Landsat-8 OLI/TIRS C1 image LC09_L2SP_173049_20221221_20221224_02_T1_SR_B: \nWhite Nile and Blue Nile confluence, Sudan (2022)", cex.main=0.9)
minor.tick(nx = 10, ny = 10, tick.ratio = 0.3)
# Plotting the histogram of the NDVI
hist(ndvi, font.main = 1, main = "NDVI values for Landsat-8 OLI/TIRS C1 image \nLC09_L2SP_173049_20221221_20221224_02_T1_SR_B: White Nile and Blue Nile confluence, Sudan (2022)", xlab = "NDVI", ylab= "Frequency",
    col = "darkgoldenrod1", xlim = c(-0.5, 1),  breaks = 30, xaxt = "n")
axis(side=1, at = seq(-0.6, 1, 0.1), labels = seq(-0.6, 1, 0.1))
minor.tick(nx = 10, ny = 10, tick.ratio = 0.3)
\end{lstlisting}

\subsection[\appendixname~\thesubsection]{R code for computing the GNDVI}
\begin{lstlisting}[language=R,caption=R code for computing the GNDVI (here: a case of 2013),style=mystyle,label={lst02}]
setwd("/Users/polinalemenkova/Documents/R/52_Image_Processing/Sudan_GNDVI")
# Green Normalized Difference Vegetation Index (GNDVI) = (NIR - Green) / (NIR + Green)
# Nir = Band 5, Green = Band 3.
vi <- function(img, k, i) {
  bk <- img[[k]]
  bi <- img[[i]]
  vi <- (bk - bi) / (bk + bi)
  return(vi)
}
# For Landsat 8-9 OLI/TIRS, NIR = 5, Green = 3.
filenames <- paste0('LC08_L2SP_173049_20131220_20200912_02_T1_SR_B', 1:7, ".tif")
filenames
landsat <- rast(filenames)
landsat
gndvi <- vi(landsat, 5, 3)
options(scipen=10000)
colors <- brewer.spectral(100)
plot(gndvi, col=colors, font.main = 1, main = "GNDVI for Landsat-8 OLI/TIRS C1 image LC08_L2SP_173049_20131220_20200912_02_T1_SR_B: \nWhite Nile and Blue Nile confluence, Sudan (2013)", cex.main=0.9)
minor.tick(nx = 10, ny = 10, tick.ratio = 0.3)
# Plotting histogram of the GNDVI
hist(gndvi, font.main = 1, main = "GNDVI values for Landsat-8 OLI/TIRS C1 image \nLC08_L2SP_173049_20131220_20200912_02_T1_SR_B: White Nile and Blue Nile confluence, Sudan (2013)", xlab = "GNDVI", ylab= "Frequency", col = "peachpuff", xlim = c(-0.5, 1),  breaks = 50, xaxt = "n")
axis(side=1, at = seq(-0.6, 1, 0.1), labels = seq(-0.6, 1, 0.1))
minor.tick(nx = 10, ny = 10, tick.ratio = 0.3)
\end{lstlisting}

\subsection[\appendixname~\thesubsection]{R code for computing the NDWI}
\begin{lstlisting}[language=R,caption=R code for computing the NDWI (here: a case of 2018),style=mystyle,label={lst03}]
vi <- function(img, k, i) {
  bk <- img[[k]]
  bi <- img[[i]]
  vi <- (bi - bk) / (bi + bk)
  return(vi)
}
# For Landsat 8-9 OLI/TIRS, NIR = 5, SWIR1 = Band 6.
filenames <- paste0('LC08_L2SP_173049_20181218_20200830_02_T1_SR_B', 1:7, ".tif")
filenames
landsat <- rast(filenames)
landsat
ndwi <- vi(landsat, 5, 6)
options(scipen=10000)
colors <- jet(100)
plot(ndwi, col=colors, font.main = 1, main = "NDWI for Landsat-8 OLI/TIRS C1 image LC08_L2SP_173049_20181218_20200830_02_T1_SR_B: \nWhite Nile and Blue Nile confluence, Sudan (2018)", cex.main=0.9)
minor.tick(nx = 10, ny = 10, tick.ratio = 0.3)
# Plotting histogram of the NDWI
hist(ndwi, font.main = 1, main = "NDWI values for Landsat-8 OLI/TIRS C1 image \nLC08_L2SP_173049_20181218_20200830_02_T1_SR_B: White Nile and Blue Nile confluence, Sudan (2018)", xlab = "NDWI", ylab= "Frequency",
    col = "mediumpurple1", xlim = c(-0.5, 1),  breaks = 50, xaxt = "n")
axis(side=1, at = seq(-0.6, 1, 0.1), labels = seq(-0.6, 1, 0.1))
minor.tick(nx = 10, ny = 10, tick.ratio = 0.3)
\end{lstlisting}

\subsection[\appendixname~\thesubsection]{R code for computing the OSAVI}
\begin{lstlisting}[language=R,caption=R code for computing the OSAVI (here: a case of 2022),style=mystyle,label={lst04}]
setwd("/Users/polinalemenkova/Documents/R/52_Image_Processing/Sudan_OSAVI")
# 1. Optimized Soil Adjusted Vegetation Index (OSAVI) = (NIR - R) / (NIR + R + 0.16), i.e., OSAVI = (Band 5 - Band 4) / (Band 5 + Band 4 + 0.16).
vi <- function(img, k, i) {
  bk <- img[[k]]
  bi <- img[[i]]
  vi <- (bk - bi) / (bk + bi + 0.16)
  return(vi)
}
# For Landsat 8-9 OLI/TIRS, NIR = 5, red = 4.
filenames <- paste0('LC09_L2SP_173049_20221221_20221224_02_T1_SR_B', 1:7, ".tif")
filenames
landsat <- rast(filenames)
landsat
osavi <- vi(landsat, 5, 4)
options(scipen=10000)
colors <- kovesi.diverging_rainbow_bgymr_45_85_c67(100)
plot(osavi, col=colors, font.main = 1, main = "OSAVI for Landsat-8 OLI/TIRS C1 image LC09_L2SP_173049_20221221_20221224_02_T1_SR_B: \nWhite Nile and Blue Nile confluence, Sudan (2022)", cex.main=0.9)
minor.tick(nx = 10, ny = 10, tick.ratio = 0.3)
# Plotting the histogram of the OSAVI
hist(osavi, font.main = 1, main = "OSAVI values for Landsat-8 OLI/TIRS C1 image \nLC09_L2SP_173049_20221221_20221224_02_T1_SR_B: White Nile and Blue Nile confluence, Sudan (2022)", xlab = "OSAVI", ylab= "Frequency", col = "dodgerblue", xlim = c(-0.5, 1),  breaks = 30, xaxt = "n")
axis(side=1, at = seq(-0.6, 1, 0.1), labels = seq(-0.6, 1, 0.1))
minor.tick(nx = 10, ny = 10, tick.ratio = 0.3)
\end{lstlisting}

\subsection[\appendixname~\thesubsection]{R code for computing the IPVI}
\begin{lstlisting}[language=R,caption=R code for computing the IPVI (here: a case of 2013),style=mystyle,label={lst05}]
library(terra)
library(RColorBrewer)
library(Hmisc)
library(pals)
#
setwd("/Users/polinalemenkova/Documents/R/52_Image_Processing/Sudan_IPVI")
# Infrared Percentage Vegetation Index (IPVI) = NIR / (NIR + Red)
# NIR = Band 5, Red = Band 4
vi <- function(img, i, k) {
  bk <- img[[k]]
  bi <- img[[i]]
  vi <- (bk) / (bi + bk)
  return(vi)
}
filenames <- paste0('LC08_L2SP_173049_20131220_20200912_02_T1_SR_B', 1:7, ".tif")
filenames
landsat <- rast(filenames)
landsat
ipvi <- vi(landsat, 5, 4)
ipvi
options(scipen=10000)
colors <- brewer.rdylgn(100)
plot(ipvi, col=colors, font.main = 1, main = "IPVI for Landsat-8 OLI/TIRS C1 image LC08_L2SP_173049_20131220_20200912_02_T1_SR_B: \nWhite Nile and Blue Nile confluence, Sudan (2013)", cex.main=0.9)
minor.tick(nx = 10, ny = 10, tick.ratio = 0.3)
# Plotting the histogram
hist(ipvi, font.main = 1, main = "IPVI values for Landsat-8 OLI/TIRS C1 image \nLC08_L2SP_173049_20131220_20200912_02_T1_SR_B: White Nile and Blue Nile confluence, Sudan (2013)", xlab = "IPVI", ylab= "Frequency",
    col = "yellowgreen", xlim = c(0, 1),  breaks = 50, xaxt = "n")
axis(side=1, at = seq(0.0, 1, 0.1), labels = seq(0.0, 1, 0.1))
minor.tick(nx = 10, ny = 10, tick.ratio = 0.3)
\end{lstlisting}

%%%%%%%%%%%%%%%%%%%%%%%%%%%%%%%%%%%%%%%%%%
\begin{adjustwidth}{-\extralength}{0cm}
%\printendnotes[custom] % Un-comment to print a list of endnotes

\reftitle{References}

%=====================================
% References, variant A: external bibliography
%=====================================
\bibliography{Sudan.bib}
%\nocite{*}

%%%%%%%%%%%%%%%%%%%%%%%%%%%%%%%%%%%%%%%%%%
\end{adjustwidth}
\end{document}
